\documentclass{article}
\usepackage[utf8]{inputenc}
\usepackage[T1]{fontenc}
\usepackage[margin=2cm]{geometry}
\usepackage{hyperref}
\sloppy
\begin{document}
\section*{Historical source}
Encyclopaedia Britannica scan, source row 1017. Parse 3.5 extracted the text. Super 3.5 VL corrected it from an enlarged crop of each block. It can still contain recognition errors.
Page headers, footers, and page numbers are removed.

\subsection*{Untitled continuation (E0074)}
\textit{Continues an entry from an earlier page.}

HEM

the sixth for the pagans, and the seventh for the hy-
pocrites of all religions.

Among Christians, there are two controverted questions in regard to hell; the one concerns locality, the other the duration of its torments. The locality of hell, and the reality of its fire, began first to be controverted by Origen. That father, interpreting the scripture account metaphorically, makes hell to consist not in external punishments, but in a consciousness or sense of guilt, and a remembrance of past pleasures. Among the moderns, Mr Whiston advanced a new hypothesis. According to him, the comets are so many hells appointed in their orbits alternately to carry the damned into the confines of the sun, there to be scorched by its violent heat, and then to return with them beyond the orb of Saturn, there to starve them in their cold and dismal regions. Another modern author, not satisfied with any hypothesis hitherto advanced, assigns the sun to be the local hell. As to the second question, viz. the duration of hell-torments, we have Origen again at the head of those who deny that they are eternal; it being that father's opinion, that not only men, but devils, after a due course of punishment suitable to their respective crimes, shall be pardoned and restored to heaven. The chief principle upon which Origen built his opinion, was the nature of punishment, which he took to be emendatory, applied only as physic for the recovery of the patient's health. The chief objection to the eternity of hell torments among modern writers, is the disproportion between temporary crimes and eternal punishments. Those who maintain the affirmative, ground their opinions on scripture accounts, which represent the pains of hell under the figure of a worm which never dies, and a fire which is not quenched; as also upon the words, "These shall go away into everlasting punishment, but the righteous into life eternal."

\subsection*{HELLEBORUS, HELLEBORE (E0075)}
HELLEBORUS, HELLEBORE, in botany, a genus of the polyandria polygynia class. It has no calix; the petals are five or more; the nectarium is tubular and bilabiated; and the capsule contains many feeds. There are five species, two of them natives of Britain, viz. the foetidus, or great bastard black hellebore; and the viridis, or wild black hellebore. The hellebore, when taken in large quantities, is poisonous; but the root, in small doses, is supposed to attenuate the humours, and to promote urinary and uterine discharges.

\subsection*{HELLENISM (E0076)}
HELLENISM, in matters of language, a phrase in the idiom, genius, or construction of the Greek tongue.

This word is only used when speaking of the authors who, writing in a different language, express themselves in a phraseology peculiar to the Greek.

\subsection*{HELLENISTIC (E0077)}
HELLENISTIC, or HELENISTIC LANGUAGE, that used by the Grecian Jews who lived in Egypt and other parts where the Greek tongue prevailed. In this language it is said the Septuagint was written, and also the books of the New Testament; and that it was thus denominated to fhew that it was Greek filled with Hebraifms and Syriacifms.

\subsection*{HELLESPONT (E0078)}
HELLESPONT, the entrance of the streights which divides Afia from Europe, and paffes from the Archi-

pelago to Constantinople. It is now called the Dardanelles, and is about two miles wide.

HELM of a ship, is a piece of timber fastened into the rudder, which comes forward into the fleerage, or place where the person at the helm steers the ship, by holding the whipstaff in his hand, which is joined to the helm. They begin however to be left off, fleering-wheels being used in their room.

There are several terms in the sea-language relating to the helm; as, bear up the helm; that is, let the ship go more large before the wind. Helm amid ship, or right the helm; that is, keep it even with the middle of the ship. Port the helm, put it over the left side of the ship. Starboard the helm, put it on the right side of the ship.

\subsection*{HELMET (E0079)}
HELMET, an ancient defensive armour worn by horsemen both in war and in tournaments. It covered both the head and face, only leaving an aperture in the front secured by bars, which was called the visor.

\subsection*{HELMINTHOLITHUS (E0080)}
HELMINTHOLITHUS, in natural history, a name given by Linnaeus to petrified bodies resembling worms. Of these he reckons four genera. 1. Petrified lithophyta, found in the mountains of Sweden. 2. Petrified shells. 3. Petrified zoophytes. 4. Petrified reptiles.

\subsection*{HELMONT (E0081)}
HELMONT, a town of the Netherlands, in the province of Dutch Brabant, situated on the river Aa : E. long. 5° 40', N. lat. 51° 30'.

\subsection*{HELMSTAT (E0082)}
HELMSTAT, a town of Germany, in the circle of Lower Saxony, and dukedom of Brunswick: E. long. 11° 15', N. lat. 52° 20'.

\subsection*{HELOTS (E0083)}
HELOTS, in Grecian antiquity, the inhabitants of Helos, a town of Laconia, conquered by the Spartans; who made them all prisoners of war, and reduced them into the condition of slaves.

The freemen of Sparta were forbidden the exercise of any mean or mechanical employment, and therefore the whole care of supplying the city with necessaries devolved upon the Helots.

HELSINGFORD a port-town of Sweden, situated on
the gulph of Finland, in 24° 6' E. long. and 60° 8'
N. lat.

\subsection*{HELSINGIA (E0084)}
HELSINGIA, a province of Sweden, bounded by the Bothnic gulf on the east, and by Dalecarlia on the west.

\subsection*{HELSINGIC CHARACTER (E0085)}
HELSINGIC CHARACTER, a peculiar kind of character, found inscribed on stones in the province of Hel- singia : the Runic and Helsingic characters may be easily transformed into each other.

\subsection*{HELSTON (E0086)}
HELSTON, a borough of Cornwall, nine miles south-
west of Falmouth: it sends two members to parlia-
ment.

\subsection*{HELVOETSLUYS (E0087)}
HELVOETSLUYS, a port-town of the united Netherlands, situated in the island of Voorn, in the province of Holland, five miles south of the Briel: it is one of the best harbours in Holland, and that to which the English packet always goes.

\subsection*{HEMEROBIUS (E0088)}
HEMEROBIUS, in zoology, a genus of insects of the neuroptera order, the characters of which are these: The mouth is furnished with two teeth; the palpi are four; the wings are deflected, but not plaited; and


\section*{Source record}
\begin{description}
\item[Scan ID] \texttt{britannica:309f765b1ecc83e0a7f7762e44b4a927a7555443:train:1017}
\item[Image file] \texttt{page-1017-5395deb0e141.jpg}
\item[Image SHA-256] \texttt{5395deb0e1413694dd6a4e16890741c049681f5838d6fc92ef14a3019a1c87a4}
\item[Dataset] NationalLibraryOfScotland/encyclopaedia\_britannica\_illustrated at revision \texttt{309f765b1ecc83e0a7f7762e44b4a927a7555443} (CC0-1.0)
\item[Extraction] model \texttt{nvidia/nemotron-3.5-parse}, raw reply \texttt{extract/raw/page-1017-5395deb0e141.json}
\item[Raw reply SHA-256] \texttt{4c778531423eb8f9d48a000a828508a0170564abd6e1f3e75d89af432d41002a}
\item[Correction] model \texttt{nvidia/Nemotron-3.5-Super-VL}, block crops, raw replies \texttt{correct/raw/crop\_nothink/page-1017-5395deb0e141-b*.json}; super\_vl: 23
\item[Running header] H E L ( 781 )
\end{description}
\end{document}
